\chapter{Complex Analysis}
\renewcommand{\currentchapterpath}{chapters/ch03/}

\begin{flushright}
\textit{“窥一斑而知全豹。” }
\footnote{Literally, ``See one spot and know the whole leopard.'' 
This condensed saying reverses the usual caution in the older idiom 
\textit{guanzhong kuibao} (viewing a leopard through a tube): 
a single observation need not be merely fragmentary when it is a 
revealing one. It expresses the ability to infer the character or 
structure of a whole from a small but telling part, provided the 
connection is understood rather than guessed.}
\end{flushright}


Complex analysis is a powerful language for physics and engineering. It combines the geometry of the complex plane with differentiation, integration, series, and residues, and it turns many real definite integrals into manageable contour integrals. This chapter follows that development from complex numbers to analytic functions, Cauchy's theory, Laurent series, and residue methods.

\chapterinput{complex_numbers}
\chapterinput{complex_functions}
\chapterinput{analytic_functions}
\chapterinput{contour_integration}
\chapterinput{complex_series}
\chapterinput{residues}
\chapterinput{applications}
\chapterinput{feynman_techniques}
\chapterinput{worked_integrals}
\chapterinput{exercises}